% =============================================================================
% CHAPTER 5: Observational Predictions & Experimental Tests
% Sources: predictions.md (PRIMARY), observational_tests.md,
%          post observational-confirmations, part2 §5.1
% =============================================================================

\chapter{Observational Predictions \& Experimental Tests}
\label{ch:predictions}

The oscillating brane theory V8.0 makes specific, testable predictions that distinguish it from standard cosmology. Three established anomalies are resolved; the definitive future test is SKA's 21\,cm reionization modulation.

% -----------------------------------------------------------------------------
\section{Timeline of Discovery}
\label{sec:timeline_discovery}

\begin{table}[h]
\centering
\caption{Observational timeline: established confirmations and future tests.}
\label{tab:discovery_timeline}
\begin{tabular}{@{}lll@{}}
\toprule
\textbf{Year} & \textbf{Instrument / Test} & \textbf{Status / Prediction} \\
\midrule
2024 & DESI detects dark energy evolution & $\checkmark$ Confirmed at $4\sigma$ \\
2025 & $S_8$ tension resolved (Yukawa) & $\checkmark$ Scale-dependent \\
2025 & Euclid first data release & $\circlearrowright$ Imminent \\
2025 & qBOUNCE / optomechanics & $\circlearrowright$ Sub-micron gravity \\
2026 & Planck CMB anomaly = ISW resonance & $\circlearrowright$ $\Delta\chi^2 = 32.9$ ($6\sigma$) \\
2027 & DESI full survey & Power spectrum modulation \\
2027+ & SKA-Low & 21\,cm modulation (\textbf{definitive}) \\
2030 & CMB-S4 & Definitive ISW signature \\
2030 & Vera Rubin/LSST & Large-scale anisotropies \\
\bottomrule
\end{tabular}
\end{table}

% -----------------------------------------------------------------------------
\section{Established Confirmations (2024--2026)}
\label{sec:confirmations}

\paragraph{Already Observed:}
\begin{itemize}
  \item \textbf{DESI 2024--2026}: Dark energy evolves with $4\sigma$ significance---exactly matching our oscillating $w(z)$ with $\phi_0 = \pi/2$.
  \item \textbf{$S_8$ tension}: Scale-dependent growth suppression via Yukawa-screened $\Geff(k)$ bridges the DES/KiDS gap.
\end{itemize}

\paragraph{Imminent Tests:}
\begin{itemize}
  \item \textbf{Euclid 2025}: Will measure $w(z)$ to 3\% precision, detecting our oscillations at $> 5\sigma$.
  \item \textbf{qBOUNCE (ILL) + levitated optomechanics}: Ultra-cold quantum neutrons and nanosphere experiments---sub-micron gravity test at $L = 0.2\,\mu$m, bypassing Casimir background.
  \item \textbf{CMB Analysis 2026}: Planck's low-$\ell$ anomaly matches our ISW resonance prediction.
\end{itemize}

% -----------------------------------------------------------------------------
\section{Dark Energy Oscillations}
\label{sec:de_oscillations}

The membrane oscillation creates a time-varying equation of state:
\begin{itemize}
  \item \textbf{Amplitude}: $A_w \ge 3 \times 10^{-3}$
  \item \textbf{Period}: $T = 2.0 \pm 0.3\Gyr$
  \item \textbf{Phase}: Maximum at $z \approx 0.5$
\end{itemize}

\textbf{Detection}: Euclid will measure $w(z)$ to 3\% precision, sufficient to detect our predicted oscillations at $> 5\sigma$ significance.

% -----------------------------------------------------------------------------
\section{ISW Resonance in the CMB}
\label{sec:isw_prediction}

The membrane oscillation creates a unique signature in the Cosmic Microwave Background through the Integrated Sachs--Wolfe effect:
\begin{itemize}
  \item \textbf{Resonance peak}: $\ell = 10$--$20$ (angular scale $\sim 12^\circ$)
  \item \textbf{Power suppression}: 16\% at resonance
  \item \textbf{Statistical significance}: $\chi^2$ improvement of 32.9 ($6\sigma$ over \LCDM)
\end{itemize}

\textbf{Detection Method}: Our $2\Gyr$ membrane oscillation ($1.6 \times 10^{-17}$~Hz) manifests through:
\begin{itemize}
  \item ISW effect: Creates resonance in CMB large-scale anisotropies at $\ell = 10$--$20$.
  \item Matter power spectrum: Periodic modulation detectable by galaxy surveys.
  \item Growth history: Time-varying structure formation rate measurable by weak lensing.
\end{itemize}

% -----------------------------------------------------------------------------
\section{Structure Growth Suppression (Scale-Dependent Yukawa Screening)}
\label{sec:s8_prediction}

The extra dimension introduces a Yukawa correction to gravity:
\begin{equation}
\Geff(k) = \GN\left(1 + \alpha\,e^{-k/\kL}\right), \quad \kL = 2\pi/L
\end{equation}

This produces $\sim 5\%$ growth suppression at non-linear scales (resolving the DES $S_8$ tension) while maintaining quasi-standard gravity at linear scales (consistent with KiDS/CMB).

% -----------------------------------------------------------------------------
\section{SKA 21\,cm Reionization Modulation --- Definitive Test}
\label{sec:ska_prediction}

The model's primary falsifiable prediction targets the 21\,cm power spectrum during the Epoch of Reionization ($6 \lesssim z \lesssim 15$). The oscillating $\Geff(k,t)$ imprints a spatial modulation on the 21\,cm brightness temperature:
\begin{equation}
\delta T_b(\vec{k}, z) \supset \Delta T_\text{osc}(k)\,\sin\!\left(\frac{2\pi\, t(z)}{T} + \phi_0\right)
\end{equation}
with characteristic amplitude $\Delta T_\text{osc} \sim 1$--$5$~mK at BAO-scale wavenumbers. \textbf{SKA-Low (2027+)} has the sensitivity and $k$-range to detect or exclude this modulation at $> 3\sigma$, constituting the definitive test.

% -----------------------------------------------------------------------------
\section{Hubble Anisotropy (Cosmicflows-4)}
\label{sec:hubble_aniso}

Spatial tension variations create directional $H_0$ differences:
\begin{equation}
\frac{\delta H}{H} \sim 10^{-3}
\end{equation}

Cosmicflows-4 bulk flow data is consistent with our elastic membrane model.

% -----------------------------------------------------------------------------
\section{Sub-Micron Gravity Tests}
\label{sec:submicron}

\begin{itemize}
  \item \textbf{qBOUNCE (ILL, Grenoble)}: Ultra-cold quantum neutrons bouncing in Earth's gravitational field map gravity at sub-micron scales. Immune to Casimir background noise that plagues mechanical force sensors.
  \item \textbf{Levitated nanoscale optomechanics}: Silica nanospheres optically trapped and cooled probe Yukawa corrections at $L = 0.2\,\mu$m directly.
\end{itemize}

Both experiments bypass the Casimir effect that blinds conventional sub-millimeter gravity tests (which operate at $\sim 1$~mm scale---completely overwhelmed by Casimir forces at our $0.2\,\mu$m target).

% -----------------------------------------------------------------------------
\section{Particle Physics Signatures}
\label{sec:particle_signatures}

\subsection{Kaluza--Klein Modes}

With $L = 0.2\,\mu$m, each Standard Model particle has an infinity of more massive copies---its excitations in the 5th dimension. The first has mass:
\begin{equation}
m_\text{KK} = \frac{\hbar}{Lc} \simeq 1\eV
\end{equation}

Too light for accelerators but potentially visible in CMB cosmology as a slight deviation in the effective number of degrees of freedom: $\Delta N_\text{eff} \sim 0.01$.

\subsection{Trans-dimensional Leakage}

Dark matter flux through the bulk induces energy ``leakage'':
\begin{equation}
\frac{\dot{\rho}}{\rho} \sim L^{-1}\,H_0 \sim 10^{-11}\text{ yr}^{-1}
\end{equation}

Future ultra-sensitive detectors (MADMAX, sub-millimeter gravity experiments) could track this slow dilution.

% -----------------------------------------------------------------------------
\section{Current Constraints}
\label{sec:current_constraints}

\begin{table}[h]
\centering
\caption{Current experimental constraints and model compatibility.}
\label{tab:constraints}
\begin{tabular}{@{}llll@{}}
\toprule
\textbf{Test} & \textbf{2024 Limit} & \textbf{Our Model} & \textbf{Verdict} \\
\midrule
Newton @ 25~$\mu$m & No deviation & $L = 0.2\,\mu$m & $\checkmark$ Invisible \\
PTA 15 years & $h_c < 3 \times 10^{-15}$ & $h_c \sim 2 \times 10^{-18}$ & $\checkmark$ Silent \\
$H_0$ dipole & $< 2\%$ & $\sim 0.01\%$ & $\checkmark$ Subtle \\
\bottomrule
\end{tabular}
\end{table}

% -----------------------------------------------------------------------------
\section{Numerical Predictions}
\label{sec:numerical_predictions}

\begin{table}[h]
\centering
\small
\caption{Quantitative predictions across observational channels.}
\label{tab:numerical_preds}
\begin{tabular}{@{}lllll@{}}
\toprule
\textbf{Observable} & \textbf{Prediction} & \textbf{Uncertainty} & \textbf{Method} & \textbf{Timeline} \\
\midrule
\multicolumn{5}{@{}l}{\textit{Fundamental Parameters}} \\
Brane tension $\tauZ$ & $7.0 \times 10^{19}$~J/m$^2$ & $\pm 15\%$ & $H_0(z)$ & Current \\
Period $T$ & $2.0\Gyr$ & $\pm 0.3\Gyr$ & ISW/BAO & Current \\
Extra dimension $L$ & $0.2\,\mu$m & Factor of 2 & KK modes & 2035+ \\
KK mass $m_\text{KK}$ & $1\eV$ & $\pm 0.5\eV$ & Cosmological & Current \\
\midrule
\multicolumn{5}{@{}l}{\textit{Cosmological Effects}} \\
$S_8$ suppression & $\sim 5\%$ (scale-dep.) & $\pm 0.5\%$ & Weak lensing & Current \\
$w(z)$ amplitude $A_w$ & 0.003 & $\pm 0.001$ & BAO + SNe & 2025+ \\
$H_0$ anisotropy & $0.01\%$ & $\pm 0.005\%$ & Precision cosmo. & 2030+ \\
\midrule
\multicolumn{5}{@{}l}{\textit{CMB ISW Effect}} \\
ISW amplitude & $\delta T/T \sim 10^{-6}$ & Factor of 2 & CMB-S4 & 2030+ \\
Angular scale & $\ell = 10$--$20$ & Full sky & Planck + CMB-S4 & Current+ \\
\midrule
\multicolumn{5}{@{}l}{\textit{Galactic Scale}} \\
MOND $a_0$ & $1.1 \times 10^{-10}$~m/s$^2$ & $\pm 5\%$ & Galaxy dynamics & Current \\
DM cross-section & $< 10^{-48}$~cm$^2$ & Lower limit & Direct detection & Current \\
\bottomrule
\end{tabular}
\end{table}

% -----------------------------------------------------------------------------
\section{Bayesian Evidence}
\label{sec:bayesian}

Current Bayesian evidence strongly favors our model:
\begin{equation}
\Delta\ln K = 3.33 \pm 0.24
\end{equation}

This represents ``strong evidence'' on the Jeffreys scale (since $e^{3.3} \approx 27$, the data make the vibrating brane scenario approximately 27 times more probable than \LCDM). The model wins because it simultaneously explains $S_8$ suppression, observed oscillation in $a(t)$, and the MOND coincidence ($a_0 \approx cH_0/2\pi$) without damaging CMB or BAO fits.

\begin{table}[h]
\centering
\caption{Jeffreys scale interpretation of Bayesian evidence.}
\begin{tabular}{@{}lll@{}}
\toprule
\textbf{$\ln K$ range} & \textbf{Interpretation} & \textbf{Our result} \\
\midrule
$< 1$ & Negligible & --- \\
$1$--$2.5$ & Modest & --- \\
$2.5$--$5$ & Strong & $\Delta\ln K = 3.33$ $\checkmark$ \\
$> 5$ & Decisive & --- \\
\bottomrule
\end{tabular}
\end{table}

This is not a definitive verdict---it is a strong signal that cosmic music might contain a real vibrato, to be confirmed (or refuted) by Euclid, DESI, and SKA in the coming years.

% -----------------------------------------------------------------------------
\section{Future Tests for 2026--2030}
\label{sec:future_tests}

\begin{table}[h]
\centering
\caption{Predictions and refutation thresholds for upcoming experiments.}
\begin{tabular}{@{}llll@{}}
\toprule
\textbf{Mission} & \textbf{Signature} & \textbf{Threshold} & \textbf{Date} \\
\midrule
Euclid & $w(z)$ sinusoidal $A \ge 3 \times 10^{-3}$ & Signal $< 5\sigma$ & 2025--27 \\
DESI Full & $\Delta P/P = 0.5\%$ at $k_0$ & Smooth spectrum & 2027 \\
CMB-S4 & ISW oscillations & No periodic pattern & 2028+ \\
SKA-Low & 21\,cm modulation 1--5~mK & No detection & 2027+ \\
H0LiCOW++ & Anisotropy $\le 0.1\%$ & Isotropy $< 0.2\%$ & 2030+ \\
\bottomrule
\end{tabular}
\end{table}
